% Figure from MATH 452 notes, section 12. Compile with the notes preamble.
\begin{tikzpicture}[scale=1.25]
  % box of half-width delta_1 = 1.6 around (x0,y0) = origin of this picture; strip |x - x0| <= delta = 1
  \fill[figB!7] (-1,-1.6) rectangle (1,1.6);
  \draw[black!55, thin] (-1.6,-1.6) rectangle (1.6,1.6);
  \draw[black!40, dashed, thin] (-1,-1.6) -- (-1,1.6);
  \draw[black!40, dashed, thin] (1,-1.6) -- (1,1.6);
  % sign conditions on top and bottom edges (Step 2)
  \draw[figR, line width=1.6pt] (-1,1.6) -- (1,1.6);
  \draw[figG, line width=1.6pt] (-1,-1.6) -- (1,-1.6);
  \node[figR, font=\scriptsize, anchor=south] at (0,1.62) {$F(x,y_0+\delta_1)>0$};
  \node[figG!80!black, font=\scriptsize, anchor=north] at (0,-1.62) {$F(x,y_0-\delta_1)<0$};
  % the level set F = 0 as the graph y = f(x)
  \draw[figB, very thick, domain=-1:1, samples=50] plot (\x, {-0.55*\x+0.18*\x*\x});
  \node[figB, font=\small, anchor=west] at (1.03,-0.33) {$y=f(x)$};
  % one vertical slice: F increases in y, crosses 0 exactly once
  \pgfmathsetmacro{\xs}{-0.55}
  \pgfmathsetmacro{\ys}{-0.55*\xs+0.18*\xs*\xs}
  \draw[black!60, thin] (\xs,-1.6) -- (\xs,1.6);
  \draw[-stealth, black!60, thin] (\xs-0.12,-1.25) -- (\xs-0.12,-0.25);
  \node[font=\scriptsize, black!70, anchor=east, align=right] at (\xs-0.14,-0.8) {$F\uparrow$\\ in $y$};
  \filldraw[figR] (\xs,\ys) circle (1.4pt);
  \node[font=\scriptsize, figR, anchor=south west, inner sep=1.5pt] at (\xs,\ys) {$(x, f(x))$};
  \filldraw[black] (0,0) circle (1.1pt);
  \node[font=\scriptsize, anchor=south west, inner sep=1.5pt] at (0.03,0.03) {$(x_0,y_0)$};
  % axis labels
  \node[font=\scriptsize, anchor=north] at (-1,-1.9) {$x_0-\delta$};
  \node[font=\scriptsize, anchor=north] at (1,-1.9) {$x_0+\delta$};
  \node[font=\scriptsize, anchor=east] at (-1.62,1.6) {$y_0+\delta_1$};
  \node[font=\scriptsize, anchor=east] at (-1.62,-1.6) {$y_0-\delta_1$};
  \node[font=\scriptsize, black!60, anchor=north west, align=left] at (1.62,1.6) {box where\\ $F_y > a/2$};
\end{tikzpicture}
